\documentclass[11pt]{article}
\usepackage[margin=1in]{geometry}
\usepackage{amsmath,amssymb,amsthm,mathtools}
\usepackage{hyperref}
\usepackage{booktabs}

\newtheorem{theorem}{Theorem}
\newtheorem{lemma}[theorem]{Lemma}
\newtheorem{proposition}[theorem]{Proposition}
\newtheorem{corollary}[theorem]{Corollary}
\theoremstyle{definition}
\newtheorem{definition}[theorem]{Definition}
\newtheorem{remark}[theorem]{Remark}

\newcommand{\Lam}{\Lambda}
\newcommand{\Z}{\mathbb{Z}}
\newcommand{\one}{\mathbf{1}}
\DeclareMathOperator{\hgt}{ht}

\title{A closed form for $[R_e(t),R_{e'}(t)]$, and the sharpness of the rigid factor}
\author{Clio}
\date{4 September 2026}

\begin{document}
\maketitle

\begin{abstract}
Let $R_e(t)$ be the operator on symmetric functions that adds an $e$-ribbon with weight
$t^{\hgt}$. Yesterday's paper proved that $(1+t)$ divides every entry of
$[R_e(t),R_{e'}(t)]$, and left two things open: whether the divisibility is sharp, and what
the quotient is. Both are settled here, by one computation. Passing to the Maya diagram, where
$R_e$ becomes a weighted bead hop $b\mapsto b+e$, the commutator is a sum over
\emph{two-hop paths}, and these fall into exactly two families. The result
(Theorem~\ref{thm:main}) is that $[R_e,R_{e'}]$ is supported on two kinds of pair
$(\lambda,\mu)$, with
\[
[R_e,R_{e'}]_{\mu\lambda}=(\alpha-\beta)\,t^{\,h-1}(1+t)
\quad\text{or}\quad
\pm\,t^{\,H-1}(1-t^2),
\]
$\alpha,\beta\in\{0,1\}$, and is zero on everything else. Hence every entry of
$Q_{e,e'}:=[R_e,R_{e'}]/(1+t)$ is $0$, $\pm t^{a}$, or $\pm t^{a}(1-t)$: at most two terms,
coefficients in $\{0,\pm1\}$. Two witness families, uniform in $e<e'$ --- $s_\varnothing\mapsto
s_{(e',1^e)}$ with entry $t^{e-1}(1+t)$, and $s_\varnothing\mapsto s_{(e',e)}$ with entry
$1-t^2$ --- show the greatest common divisor of the entries is \emph{exactly} $t+1$
(Corollary~\ref{cor:sharp}), which was the open half of the programme's rigidity statement.
The geometric content is a commutation criterion: two ribbon additions commute, \emph{as graded
operators}, unless the intervals of diagonals they occupy properly cross
(Corollary~\ref{cor:diag}). Finally, a bounded prior-art phase settles the status of yesterday's
Theorem~A: it is \textbf{classical}, being --- up to the involution $\omega$ and a global sign,
a dictionary I verified in my own engine on two independent routes --- the $q$-Murnaghan--Nakayama
rule for Iwahori--Hecke algebra characters (\S\ref{sec:prior}).
\end{abstract}

\section{What was open}

The standing result is \texttt{proofs/2026-09-03-Q75-multiplication-defect.tex}. Recall the
object. For $0\le h\le e-1$ let $N_e^{(h)}$ be the operator on the ring $\Lam$ of symmetric
functions over $\Z[t]$ with $N_e^{(h)}s_\lambda=\sum_\mu s_\mu$, summed over those $\mu$
obtained from $\lambda$ by \textbf{adding} an $e$-ribbon spanning exactly $h+1$ rows, and put
\[
R_e(t)=\sum_{h=0}^{e-1}t^hN_e^{(h)},
\qquad\text{so}\qquad
R_e(t)_{\mu\lambda}=\begin{cases}t^{h(\mu/\lambda)}&\mu/\lambda\text{ an $e$-ribbon},\\0&\text{else.}\end{cases}
\]
Here an $e$-ribbon (rim hook, border strip) is a connected skew shape of size $e$ with no
$2\times2$ square, and its height $h$ is one less than its number of rows.

Yesterday proved $R_e(-1)=M_{p_e}$, multiplication by the power sum, and deduced from the
commutativity of $\Lam$ that $(1+t)$ divides every entry of $[R_e(t),R_{e'}(t)]$. Two questions
were left open, and are the subject of this paper.

\begin{itemize}
\item[(S)] \emph{Sharpness.} Is the greatest common divisor of the entries of
$[R_e,R_{e'}]$ exactly $t+1$, rather than a higher power? Equivalently, is
$Q_{e,e'}:=[R_e,R_{e'}]/(1+t)$ nonzero at $t=-1$? Observed on three pairs $(e,e')$; not proved.
\item[(C)] \emph{Closed form.} What is $Q_{e,e'}$?
\end{itemize}

Both are answered by Theorem~\ref{thm:main}. The move that makes them finite is to stop working
with skew shapes and work with beads.

\section{The bead model}\label{sec:abacus}

\begin{definition}\label{def:maya}
For a partition $\lambda$ put $B(\lambda)=\{\lambda_i-i:i\ge1\}\subset\Z$ (with $\lambda_i=0$ for
$i>\ell(\lambda)$). This is the \emph{Maya diagram} of $\lambda$: it contains every sufficiently
negative integer and omits every sufficiently large one. Writing its elements in decreasing order
$b_1>b_2>\cdots$, one recovers $\lambda_i=b_i+i$, so $\lambda\mapsto B(\lambda)$ is injective.
\end{definition}

The following is classical (it is the abacus of James--Kerber); I give the proof because the
programme's height convention is the part of the object that is mine, and because the last
fortnight's failures were all in the evidence layer. Nothing below is quoted from memory.

\begin{proposition}[bead hops]\label{prop:hop}
Let $\lambda$ be a partition, $B=B(\lambda)$, and $e\ge1$.
\begin{enumerate}
\item[(i)] The partitions $\mu$ such that $\mu/\lambda$ is an $e$-ribbon are exactly the
partitions with $B(\mu)=(B\setminus\{b\})\cup\{b+e\}$ for some $b\in B$ with $b+e\notin B$; and
$b$ is determined by $\mu$, being the unique element of $B(\lambda)\setminus B(\mu)$.
\item[(ii)] For such a $\mu$, the height is $h(\mu/\lambda)=\#\bigl(B\cap(b,b+e)\bigr)$, the
number of beads strictly jumped over.
\item[(iii)] The cells of $\mu/\lambda$ have contents (i.e.\ values of $j-i$ on a cell in row $i$,
column $j$) exactly $b+1,b+2,\dots,b+e$, each once.
\item[(iv)] Listing the cells of $\mu/\lambda$ as $c_1,\dots,c_e$ in increasing content --- which
is the order along the ribbon from its bottom-left end --- the cell $c_{j+1}$ lies directly above
$c_j$ if $b+j\in B$, and directly to the right of $c_j$ if $b+j\notin B$.
\end{enumerate}
\end{proposition}

\begin{proof}
Let $b\in B$ with $b+e\notin B$, let $k=\#\bigl(B\cap(b,b+e)\bigr)$, and let $b=b_r$ in the
decreasing enumeration of $B$. The elements of $B$ lying in $(b,b+e)$ are then exactly
$b_{r-1},b_{r-2},\dots,b_{r-k}$, and $b_{r-k-1}>b+e$ (using $b+e\notin B$). Set
$B'=(B\setminus\{b\})\cup\{b+e\}$. Its decreasing enumeration is
\[
b_1,\ \dots,\ b_{r-k-1},\ \underbrace{b+e}_{\text{position }r-k},\ b_{r-k},\ b_{r-k+1},\ \dots,\ b_{r-1},\ b_{r+1},\ \dots,
\]
so, writing $\mu$ for the partition with $B(\mu)=B'$ and using $\mu_i=b'_i+i$,
\begin{equation}\label{eq:rows}
\mu_i=\lambda_i\ (i<r-k\text{ or }i>r),\qquad
\mu_{r-k}=b+e+r-k,\qquad
\mu_i=\lambda_{i-1}+1\ (r-k<i\le r).
\end{equation}
In particular $\mu\supseteq\lambda$: for $r-k<i\le r$ we have $\mu_i=\lambda_{i-1}+1>\lambda_i$,
and $\mu_{r-k}=b+e+r-k>b_{r-k}+r-k=\lambda_{r-k}$ because $b_{r-k}<b+e$.

Now read off $\mu/\lambda$ from \eqref{eq:rows}. It is empty outside rows $r-k,\dots,r$. In row
$i$ with $r-k<i\le r$ it occupies the columns $\lambda_i+1,\dots,\lambda_{i-1}+1$, which is
$\lambda_{i-1}+1-\lambda_i=b_{i-1}-b_i\ge1$ cells; in row $r-k$ it occupies the columns
$\lambda_{r-k}+1,\dots,b+e+r-k$, which is $b+e-b_{r-k}\ge1$ cells. Summing,
\[
|\mu/\lambda|=(b+e-b_{r-k})+\sum_{i=r-k+1}^{r}(b_{i-1}-b_i)=(b+e-b_{r-k})+(b_{r-k}-b_r)=e,
\]
since $b_r=b$. (When $k=0$ the middle sum is empty and the single row $r$ has $b+e-b_r=e$ cells.)
Rows $i-1$ and $i$, for $r-k<i\le r$, share exactly the column $\lambda_{i-1}+1$ --- it is the
last column of row $i$ and the first of row $i-1$ --- and rows $i$ and $i-2$ share none, since
row $i$ stops at column $\lambda_{i-1}+1$ and row $i-2$ starts at column
$\lambda_{i-2}+1>\lambda_{i-1}+1$. A skew shape whose nonempty rows are consecutive and whose
consecutive rows overlap in exactly one column is connected and contains no $2\times2$ square;
so $\mu/\lambda$ is an $e$-ribbon, and it occupies $k+1$ rows, giving (ii).

For (iii): the contents occupied in row $i$ ($r-k<i\le r$) are
$(\lambda_i+1)-i,\dots,(\lambda_{i-1}+1)-i$, i.e.\ the integers in $(b_i,b_{i-1}]$; and in row
$r-k$ they are the integers in $(b_{r-k},b+e]$. These intervals are consecutive and disjoint, and
their union is $(b_r,b+e]=(b,b+e]$. This proves (iii), and shows that increasing content is the
same as travelling along the ribbon from row $r$ to row $r-k$, i.e.\ from the bottom-left end.
For (iv): within a row the content increases by $1$ as the column does, so $c_{j+1}$ is to the
right of $c_j$ unless $c_j$ is the last cell of its row. By the description just given, the last
cell of row $i$ has content $b_{i-1}$, and the first cell of row $i-1$ --- which is the next one
in content order --- lies in column $\lambda_{i-1}+1$, the same column. So the ribbon steps up
between contents $d$ and $d+1$ exactly when $d\in\{b_{r-1},\dots,b_{r-k}\}=B\cap(b,b+e)$. Since
$c_j$ has content $b+j$, this is the statement of (iv).

Finally, the converse half of (i). Let $\mu/\lambda$ be an $e$-ribbon, occupying rows
$r-k,\dots,r$ for some $k\ge0$. Consecutive rows of a ribbon overlap in exactly one column,
which forces $\mu_i=\lambda_{i-1}+1$ for $r-k<i\le r$, and $\mu_i=\lambda_i$ off
$[r-k,r]$. Hence $b'_i=b_{i-1}$ for $r-k<i\le r$ and $b'_i=b_i$ for $i<r-k$ and $i>r$, so
$B(\mu)$ and $B(\lambda)$ differ in exactly one element each: $B(\lambda)\setminus B(\mu)=\{b_r\}$
and $B(\mu)\setminus B(\lambda)=\{b'_{r-k}\}$. Comparing sizes, $|\mu|-|\lambda|=b'_{r-k}-b_r=e$.
Setting $b=b_r$ gives $B(\mu)=(B\setminus\{b\})\cup\{b+e\}$, and $b+e\notin B$ because
$b+e=b'_{r-k}\notin B(\lambda)$.
\end{proof}

\begin{corollary}\label{cor:Rmaya}
Identify $s_\lambda$ with $B(\lambda)$. Then
\[
R_e(t)\,\lvert B\rangle=\sum_{\substack{b\in B\\ b+e\notin B}}
t^{\,\#(B\cap(b,b+e))}\ \bigl\lvert (B\setminus\{b\})\cup\{b+e\}\bigr\rangle .
\]
\end{corollary}

Everything from here on is a statement about beads hopping to the right on $\Z$.

\section{Two-hop paths}

\begin{definition}
A \emph{hop} is a pair $(u,d)$ with $u\in\Z$ and $d\ge1$. It is \emph{applicable} to a Maya
diagram $B$ if $u\in B$ and $u+d\notin B$, and it then produces $B\cdot(u,d)=(B\setminus\{u\})\cup\{u+d\}$
with \emph{weight} $t^{\#(B\cap(u,u+d))}$.

Fix $e\ne e'$ and Maya diagrams $B,B'$. A \emph{path} from $B$ to $B'$ is a sequence of two hops
$(u_1,d_1),(u_2,d_2)$ with $(d_1,d_2)\in\{(e',e),(e,e')\}$, the first applicable to $B$, the
second to $B_1:=B\cdot(u_1,d_1)$, and $B_1\cdot(u_2,d_2)=B'$. Its \emph{sign} is $+1$ if
$(d_1,d_2)=(e',e)$ and $-1$ if $(d_1,d_2)=(e,e')$, and its weight is the product of the two hop
weights.
\end{definition}

By Corollary~\ref{cor:Rmaya}, expanding $R_eR_{e'}-R_{e'}R_e$ (the rightmost operator acts first)
gives at once:

\begin{lemma}\label{lem:sum}
$[R_e,R_{e'}]_{\mu\lambda}=\sum_{P}\operatorname{sign}(P)\cdot\operatorname{wt}(P)$, over all
paths $P$ from $B(\lambda)$ to $B(\mu)$.
\end{lemma}

\begin{lemma}[trichotomy]\label{lem:tri}
Let $P=\bigl((u_1,d_1),(u_2,d_2)\bigr)$ be a path from $B$ to $B'$, and put $v_i=u_i+d_i$. Exactly
one of the following holds.
\begin{enumerate}
\item[\rm(F)] $u_2=v_1$ (the same bead hops twice). Then $B\setminus B'=\{u_1\}$ and
$B'\setminus B=\{u_1+e+e'\}$.
\item[\rm(B)] $v_2=u_1$ (the second bead lands on the site the first vacated). Then
$B\setminus B'=\{u_2\}$ and $B'\setminus B=\{u_2+e+e'\}$.
\item[\rm(S)] neither. Then $u_1\ne u_2$, $v_1\ne v_2$, $B\setminus B'=\{u_1,u_2\}$ and
$B'\setminus B=\{v_1,v_2\}$.
\end{enumerate}
\end{lemma}

\begin{proof}
$B_1=(B\setminus\{u_1\})\cup\{v_1\}$ and $B'=(B_1\setminus\{u_2\})\cup\{v_2\}$. Applicability of
the second hop says $u_2\in B_1$ and $v_2\notin B_1$. If $u_2=v_1$ we are in (F), and then
$B'=(B\setminus\{u_1\})\cup\{v_2\}$ with $v_2=u_1+d_1+d_2=u_1+e+e'$; note $v_2\ne u_1$ since
$e+e'>0$. Otherwise $u_2\in B$ and $u_2\ne u_1$. If moreover $v_2=u_1$ we are in (B), and
$B'=(B\setminus\{u_2\})\cup\{v_1\}$ with $v_1=u_1+d_1=v_2+d_1=u_2+d_2+d_1=u_2+e+e'$. Otherwise
$v_2\notin B$ and $v_2\ne v_1$, giving (S). Cases (F) and (B) are exclusive: they would give
$v_2=u_1<v_1=u_2<v_2$.
\end{proof}

So the support of $[R_e,R_{e'}]$ is confined to pairs $(\lambda,\mu)$ with
$\#\bigl(B(\lambda)\setminus B(\mu)\bigr)\in\{1,2\}$, and the two values are handled separately.

\subsection*{The ribbon case}

\begin{lemma}\label{lem:type2}
Suppose $B\setminus B'=\{b\}$. Then $B'\setminus B=\{b+e+e'\}$ (else there is no path), and,
writing $N=\#\bigl(B\cap(b,b+e+e')\bigr)$, $\alpha=\one[b+e\in B]$, $\beta=\one[b+e'\in B]$,
\[
\sum_{P}\operatorname{sign}(P)\operatorname{wt}(P)=(\alpha-\beta)\,t^{\,N-1}(1+t).
\]
Moreover $N\ge1$ whenever $\alpha\ne\beta$, so the exponent is non-negative.
\end{lemma}

\begin{proof}
By Lemma~\ref{lem:tri} every path is of type (F) or (B), and in both cases the vacated site is
$b$ and the filled site is $b+e+e'$. There are four candidates, indexed by the type and by
$(d_1,d_2)\in\{(e',e),(e,e')\}$.

\emph{Type (F), $(d_1,d_2)$.} The hops are $(b,d_1)$ then $(b+d_1,d_2)$. Applicability: $b\in B$
holds; $b+d_1\notin B$ is a genuine condition; then $b+d_1\in B_1$ holds and
$b+e+e'\notin B_1$ holds because $b+e+e'\notin B$ and $b+e+e'\ne b+d_1$. So this path exists iff
$b+d_1\notin B$. Its weight is $t^{\#(B\cap(b,b+d_1))}\cdot t^{\#(B_1\cap(b+d_1,b+e+e'))}$, and
$B_1=(B\setminus\{b\})\cup\{b+d_1\}$ agrees with $B$ on the open interval $(b+d_1,b+e+e')$;
hence the weight is
$t^{\#(B\cap(b,b+d_1))+\#(B\cap(b+d_1,b+e+e'))}=t^{N-\one[b+d_1\in B]}=t^{N}$.

\emph{Type (B), $(d_1,d_2)$.} Here $u_2=b$, $v_2=b+d_2$, $u_1=v_2=b+d_2$ and $v_1=b+e+e'$. The
hops are $(b+d_2,d_1)$ then $(b,d_2)$. Applicability: $b+d_2\in B$ is a genuine condition;
$b+e+e'\notin B$ holds; then $b\in B_1$ holds and $b+d_2\notin B_1$ holds since $b+d_2$ was
removed and $b+d_2\ne b+e+e'$. So this path exists iff $b+d_2\in B$. Its weight is
$t^{\#(B\cap(b+d_2,b+e+e'))}\cdot t^{\#(B_1\cap(b,b+d_2))}$, and $B_1$ agrees with $B$ on
$(b,b+d_2)$; hence the weight is $t^{N-\one[b+d_2\in B]}=t^{N-1}$.

Collecting, with sign $+$ for $(d_1,d_2)=(e',e)$ and $-$ for $(e,e')$:
\[
\bigl((1-\beta)t^{N}+\alpha t^{N-1}\bigr)-\bigl((1-\alpha)t^{N}+\beta t^{N-1}\bigr)
=(\alpha-\beta)\bigl(t^{N}+t^{N-1}\bigr)=(\alpha-\beta)t^{N-1}(1+t).
\]
If $\alpha\ne\beta$ then one of $b+e,b+e'$ lies in $B$, and both lie strictly inside
$(b,b+e+e')$ because $e,e'\ge1$; so $N\ge1$.
\end{proof}

\subsection*{The two-bead case}

\begin{lemma}[the assignment is forced]\label{lem:assign}
Let $e\ne e'$ and suppose $B\setminus B'=\{p,q\}$ and $B'\setminus B=\{x,y\}$ with $p\ne q$,
$x\ne y$. There is at most one bijection $\phi:\{p,q\}\to\{x,y\}$ whose displacement multiset
$\{\phi(p)-p,\ \phi(q)-q\}$ equals $\{e,e'\}$.
\end{lemma}

\begin{proof}
Suppose both bijections work: $\{x-p,\,y-q\}=\{e,e'\}$ and $\{y-p,\,x-q\}=\{e,e'\}$. If
$x-p=e$ then $y-p\ne e$ (else $x=y$), so $y-p=e'$; since $e\ne e'$ the second condition forces
$x-q=e$, whence $p=q$. The case $x-p=e'$ is symmetric.
\end{proof}

\begin{lemma}\label{lem:type1}
Let $e\ne e'$, suppose $B\setminus B'$ has two elements, and suppose the bijection of
Lemma~\ref{lem:assign} exists; write $b$ for the bead displaced by $e$ and $c$ for the bead
displaced by $e'$, so $B'=(B\setminus\{b,c\})\cup\{b+e,c+e'\}$. Put
$H=\#\bigl(B\cap(b,b+e)\bigr)+\#\bigl(B\cap(c,c+e')\bigr)$. Then
\[
\sum_{P}\operatorname{sign}(P)\operatorname{wt}(P)=
\begin{cases}
+\,t^{\,H-1}(1-t^2)&\text{if }b<c<b+e<c+e',\\
-\,t^{\,H-1}(1-t^2)&\text{if }c<b<c+e'<b+e,\\
0&\text{otherwise,}
\end{cases}
\]
and $H\ge1$ in the first two cases. If the bijection does not exist, the sum is $0$.
\end{lemma}

\begin{proof}
By Lemmas~\ref{lem:tri} and \ref{lem:assign} there are at most two paths, differing only in the
order of the two hops $(b,e)$ and $(c,e')$. Both exist. Indeed $b,c\in B$ and
$b+e,c+e'\notin B$ (they lie in $B'\setminus B$), and $b+e\ne c+e'$; for the order
``$c$ first'' one needs $c+e'\notin B$ and then $b+e\notin B_c$ where
$B_c=(B\setminus\{c\})\cup\{c+e'\}$, which holds since $b+e\notin B$ and $b+e\ne c+e'$; the other
order is symmetric. The order ``$c$ first, then $b$'' has $(d_1,d_2)=(e',e)$ and so sign $+$.

Weights. With $B_c$ as above and $B_b=(B\setminus\{b\})\cup\{b+e\}$,
\[
\#\bigl(B_c\cap(b,b+e)\bigr)=\#\bigl(B\cap(b,b+e)\bigr)-\delta_1+\delta_2,\qquad
\#\bigl(B_b\cap(c,c+e')\bigr)=\#\bigl(B\cap(c,c+e')\bigr)-\delta_3+\delta_4,
\]
where $\delta_1=\one[c\in(b,b+e)]$, $\delta_2=\one[c+e'\in(b,b+e)]$, $\delta_3=\one[b\in(c,c+e')]$
and $\delta_4=\one[b+e\in(c,c+e')]$. Hence the total is
\begin{equation}\label{eq:delta}
t^{H}\bigl(t^{\delta_2-\delta_1}-t^{\delta_4-\delta_3}\bigr).
\end{equation}
Note $b\ne c$, and $b+e\ne c$ and $c+e'\ne b$ because $B\setminus B'$ and $B'\setminus B$ are
disjoint, and $b+e\ne c+e'$.

Assume first $b<c$; then $\delta_3=0$, $\delta_1=\one[c<b+e]$,
$\delta_2=\one[c+e'<b+e]$ (as $c+e'>c>b$) and $\delta_4=\one[c<b+e<c+e']$.
\begin{itemize}
\item If $c>b+e$: then $\delta_1=0$; also $c+e'>c>b+e$ so $\delta_2=0$, and $\delta_4=0$. By
\eqref{eq:delta} the total is $t^H(1-1)=0$.
\item If $c<b+e$: then $\delta_1=1$ and $\delta_4=\one[b+e<c+e']$. Since $b+e\ne c+e'$, either
$c+e'<b+e$, giving $\delta_2=1,\delta_4=0$ and total $t^H(t^0-t^0)=0$; or $c+e'>b+e$, giving
$\delta_2=0,\delta_4=1$ and total $t^H(t^{-1}-t)=t^{H-1}(1-t^2)$. The second alternative is
precisely $b<c<b+e<c+e'$, and there $c\in B\cap(b,b+e)$, so $H\ge1$.
\end{itemize}
The case $c<b$ is obtained by exchanging the roles of $(b,e)$ and $(c,e')$, which exchanges
$(\delta_1,\delta_2)$ with $(\delta_3,\delta_4)$ and therefore negates \eqref{eq:delta}.
\end{proof}

\section{The closed form}

\begin{theorem}\label{thm:main}
Let $e\ne e'$ and let $\lambda,\mu$ be partitions; write $B=B(\lambda)$, $B'=B(\mu)$. Then
$[R_e(t),R_{e'}(t)]_{\mu\lambda}$ is:
\begin{enumerate}
\item[\rm(i)] if $B\setminus B'=\{b\}$ and $B'\setminus B=\{b+e+e'\}$ --- equivalently, if
$\mu/\lambda$ is an $(e+e')$-ribbon, in which case $h:=h(\mu/\lambda)=\#(B\cap(b,b+e+e'))$ ---
\[
\bigl(\one[b+e\in B]-\one[b+e'\in B]\bigr)\ t^{\,h-1}\,(1+t);
\]
\item[\rm(ii)] if $B\setminus B'=\{b,c\}$ and $B'\setminus B=\{b+e,c+e'\}$ (an assignment that is
unique if it exists), with $H=\#(B\cap(b,b+e))+\#(B\cap(c,c+e'))$ ---
\[
\begin{cases}
+\,t^{\,H-1}(1-t^2)&\text{if }b<c<b+e<c+e',\\
-\,t^{\,H-1}(1-t^2)&\text{if }c<b<c+e'<b+e,\\
0&\text{otherwise;}
\end{cases}
\]
\item[\rm(iii)] $0$ in every other case.
\end{enumerate}
All exponents occurring with a nonzero coefficient are non-negative.
\end{theorem}

\begin{proof}
Lemma~\ref{lem:sum} reduces the entry to a signed sum over paths; Lemma~\ref{lem:tri} splits the
paths into the case $\#(B\setminus B')=1$, handled by Lemma~\ref{lem:type2}, and the case
$\#(B\setminus B')=2$, handled by Lemmas~\ref{lem:assign} and \ref{lem:type1}; and no other value
of $\#(B\setminus B')$ admits a path. In case (i) the description of $\mu/\lambda$ as an
$(e+e')$-ribbon, and the identification of $h$, are Proposition~\ref{prop:hop}(i)--(ii) applied
with ribbon size $e+e'$.
\end{proof}

\begin{corollary}[the shape of the quotient]\label{cor:shape}
Every entry of $Q_{e,e'}:=[R_e,R_{e'}]/(1+t)$ is $0$, or $\pm t^{a}$, or $\pm t^{a}(1-t)$ for
some $a\ge0$. In particular every entry has at most two terms and all its coefficients lie in
$\{0,\pm1\}$; the monomial entries are exactly those on $(e+e')$-ribbons.
\end{corollary}

\begin{proof}
Divide the two cases of Theorem~\ref{thm:main} by $1+t$, using $1-t^2=(1-t)(1+t)$.
\end{proof}

This also re-proves, without reference to $\Lam$ being commutative, that $(1+t)$ divides
$[R_e,R_{e'}]$: both surviving shapes carry the factor visibly.

\section{Sharpness}

\begin{lemma}[two witness families]\label{lem:wit}
Let $2\le e<e'$. Then
\[
[R_e,R_{e'}]_{(e',1^e),\,\varnothing}=t^{\,e-1}(1+t),
\qquad
[R_e,R_{e'}]_{(e',e),\,\varnothing}=1-t^2 .
\]
\end{lemma}

\begin{proof}
$B(\varnothing)=\Z_{<0}$, so every negative integer is a bead and no non-negative integer is.

\emph{First family.} Take $b=-1-e$ and hop it to $b+e+e'=e'-1$. The beads strictly between are
$-e,\dots,-1$, so $h=e$. Here $b+e=-1\in B$ and $b+e'=e'-e-1\ge0$ is not in $B$, so
$\one[b+e\in B]-\one[b+e'\in B]=1$ and Theorem~\ref{thm:main}(i) gives $t^{e-1}(1+t)$. The
target: $B'=\bigl(\Z_{<0}\setminus\{-1-e\}\bigr)\cup\{e'-1\}$ has decreasing enumeration
$e'-1,-1,-2,\dots,-e,-e-2,-e-3,\dots$, so by $\mu_i=b'_i+i$ we get
$\mu_1=e'$, $\mu_2=\dots=\mu_{e+1}=1$ and $\mu_i=0$ beyond; that is $\mu=(e',1^e)$.

\emph{Second family.} Take $b=-2$ hopping by $e$ and $c=-1$ hopping by $e'$. Both are beads;
$b+e=e-2\ge0$ and $c+e'=e'-1\ge0$ are not, and they are distinct. The interval $(b,b+e)=(-2,e-2)$
meets $B$ only in $\{-1\}$ (as $e\ge2$), and $(c,c+e')=(-1,e'-1)$ meets $B$ not at all; so $H=1$.
The crossing condition $b<c<b+e<c+e'$ reads $-2<-1<e-2<e'-1$, which holds because $e\ge2$ and
$e<e'$. Theorem~\ref{thm:main}(ii) gives $t^{0}(1-t^2)=1-t^2$. The target:
$B'=\bigl(\Z_{<0}\setminus\{-2,-1\}\bigr)\cup\{e-2,e'-1\}$ enumerates as
$e'-1,e-2,-3,-4,\dots$, giving $\mu_1=e'$, $\mu_2=e$, $\mu_i=0$ beyond; that is $\mu=(e',e)$.
\end{proof}

\begin{corollary}[sharpness; question (S)]\label{cor:sharp}
For all $2\le e<e'$ the greatest common divisor of the entries of $[R_e(t),R_{e'}(t)]$ is exactly
$t+1$. Equivalently $Q_{e,e'}(-1)\ne0$; indeed $Q_{e,e'}(-1)_{(e',1^e),\varnothing}=(-1)^{e-1}$.
\end{corollary}

\begin{proof}
Divisibility by $t+1$ is Corollary~\ref{cor:shape} (or Cor.~T3 of the previous paper), so the gcd
is $(t+1)g$ where $g$ is the gcd of the entries of $Q_{e,e'}$. By Lemma~\ref{lem:wit},
$Q_{e,e'}$ has an entry $t^{e-1}$ and an entry $1-t$. Hence $g$ divides $t^{e-1}$, so $g$ is a
power of $t$ up to a unit; and $g$ divides $1-t$, which is not divisible by $t$. So $g$ is a unit
and the gcd is $t+1$. Evaluating $t^{e-1}$ at $t=-1$ gives $(-1)^{e-1}\ne0$.
\end{proof}

Both families are uniform in $(e,e')$, which was the form the brief asked for; the weaker
fallback (a non-uniform witness, or only $e'=e+1$) is not needed.

\section{Reading the answer on Young diagrams}

Theorem~\ref{thm:main} is stated on beads because that is where it is finite. Proposition
\ref{prop:hop}(iii)--(iv) translates it back.

\begin{corollary}[the ribbon entries]\label{cor:ribbon}
Let $\mu/\lambda$ be an $(e+e')$-ribbon of height $h$, and list its cells $c_1,\dots,c_{e+e'}$
from its bottom-left end (equivalently, in increasing content). For $1\le j<e+e'$ set
$\epsilon_j=1$ if $c_{j+1}$ lies directly above $c_j$ and $\epsilon_j=0$ if it lies directly to
its right. Then
\[
[R_e,R_{e'}]_{\mu\lambda}=(\epsilon_e-\epsilon_{e'})\,t^{\,h-1}(1+t),
\]
which is nonzero exactly when the ribbon turns upward at exactly one of its two cut points.
\end{corollary}

\begin{proof}
Immediate from Theorem~\ref{thm:main}(i) and Proposition~\ref{prop:hop}(iv), which says
$\epsilon_j=\one[b+j\in B]$.
\end{proof}

The first witness family is the smallest instance: $\mu/\varnothing=(e',1^e)$ is a hook, its
cells read from the foot of the leg, and it turns upward at cut point $e$ (the leg-to-arm corner)
and rightward at cut point $e'$ (inside the arm, since $e'>e$).

\begin{corollary}[a commutation criterion]\label{cor:diag}
By Proposition~\ref{prop:hop}(iii), the $e$-ribbon added by the hop $b\mapsto b+e$ occupies
exactly the diagonals (contents) in the interval $(b,b+e]$. Case (ii) of Theorem~\ref{thm:main}
therefore says: two ribbon additions, of sizes $e$ and $e'$ and by distinct beads, contribute
equally to the two orders --- that is, they commute \emph{as graded operators} --- unless the two
intervals of diagonals they occupy \emph{properly cross}, meaning they overlap and neither
contains the other. Disjoint intervals commute, and so do nested ones.
\end{corollary}

The nested case is the one worth noticing, because it is the case a parity or degree argument
would be least likely to predict: a small ribbon sitting entirely inside the diagonal span of a
large one still commutes with it. Specialised to $e=e'=n$, the criterion reads $|b-c|\ge n$,
which is the criterion recorded in the accepted answer to MathOverflow~292312 (Alexandersson's
question, answered by Ziqing Xiang) for when two rim-hook removals may be swapped. I have
\emph{not} read that thread at source this session --- it is recorded in my browse log of
2026-09-04 at the level of an agent's check --- so I record the agreement as a consistency
observation, not as a citation.

\begin{remark}
Two structural facts fall out. First, $[R_e,R_{e'}]$ vanishes at $t=1$ on every two-bead entry
(where the factor is $1-t^2$) and equals $\pm2$ on the ribbon entries; so at $t=1$ the commutator
is supported on $(e+e')$-ribbons alone. Second, $(1+t)^2$ never divides an entry: the ribbon
entries have $Q$-value $\pm t^{h-1}$, and the two-bead entries $\pm t^{H-1}(1-t)$, neither of
which vanishes at $t=-1$. That is a stronger statement than Corollary~\ref{cor:sharp}: not merely
some entry, but \emph{every} nonzero entry of $[R_e,R_{e'}]$ is divisible by $1+t$ exactly once.
\end{remark}

\section{Prior art}\label{sec:prior}

\subsection*{Theorem A of the previous paper is classical}

The 2026-09-03 paper proved (its Theorem~A) that for a skew shape $S$ of size $e$,
$\langle s_S,f_e(t)\rangle=t^{h(S)}(1+t)^{c(S)-1}$ if $S$ has no $2\times2$ square and $0$
otherwise, where $f_e(t)=\sum_{k}t^ks_{(e-k,1^k)}$, and left its prior-art status
``unsettled, flagged''. A bounded search this session settles it: \textbf{it is classical}.

Jing and Liu, \emph{Murnaghan--Nakayama rule\ldots}, arXiv \texttt{2310.15730}, define in \S4.2
\[
(a-1)\sum_{r\ge0}\tilde h_r(a)z^r=\prod_i\frac{1-x_iz}{1-ax_iz},
\qquad
wt(\theta;t)=(t-1)^{m-1}\prod_{i=1}^{m}(-1)^{r(\xi_i)-1}t^{c(\xi_i)-1}
\]
($\theta$ a generalised border strip with components $\xi_1,\dots,\xi_m$, of $r(\xi_i)$ rows and
$c(\xi_i)$ columns), and state, as a corollary, $\tilde h_r(q)s_\mu=\sum_\lambda
wt(\lambda/\mu;q)s_\lambda$ summed over $\lambda$ with $\lambda/\mu$ an $r$-generalised border
strip. I read all three displays myself in the \texttt{2310.15730} e-print source
(\texttt{mn.tex}, lines 626, 728, 744), in context, not from a quotation.

Their attributions, stated precisely because the difference matters. The corollary is introduced
as ``The following Murnaghan--Nakayama rule (see [Ram], [Hal], [JL2]) for $H_n(q)$'' --- i.e.\
credited to Ram, \emph{Invent.\ Math.} \textbf{106} (1991), Halverson, \emph{JCTA} \textbf{71}
(1995), and an earlier Jing--Liu paper. The weight lemma \texttt{t:wt1} is introduced as
``similar to [HR, (6.7)]'' --- \emph{similar to}, not \emph{due to} --- referring to
Halverson--Ram, \emph{Adv.\ Math.} \textbf{116} (1995), and Jing and Liu ``furnish a proof for
completeness''. So the claim ``classical'' rests on the Ram/Halverson attribution of the
\emph{rule}, which I have not followed to its source.

One thing read at source that bears on the previous paper's other gap: Jing and Liu describe
Ram's Frobenius formula as expressing Hecke-algebra characters ``in terms of one-row
Hall--Littlewood functions and Schur symmetric functions''. That is independent corroboration,
from a document I have read, that $\tilde h_r$ --- and hence $f_e(t)$ --- is a one-row
Hall--Littlewood function.

I checked the dictionary myself, in my own engine, by two routes that share no code, rather than
accepting the algebra second-hand:
\begin{itemize}
\item \emph{Generating functions.} Their right-hand side at $a=-t$ is $E(-z)H(-tz)$; mine
(Lemma~gf of the previous paper) is $\sum_e(1+t)f_e(t)z^e=H(z)E(tz)$, and applying $\omega$ and
$z\mapsto-z$ carries one to the other. Coefficient by coefficient this predicts
$\tilde h_e(-t)=(-1)^{e-1}\omega f_e(t)$, confirmed in the Schur basis for $e\le5$.
\item \emph{Weights.} Using $r(\xi)+c(\xi)=|\xi|+1$, their $wt(\theta;-t)$ should equal
$(-1)^{|\theta|-1}$ times my weight for the \emph{conjugate} shape $\theta'$. Confirmed on all
$541$ skew shapes $\mu/\lambda$ with $|\lambda|\le5$ and $|\mu/\lambda|\le4$ ($527$ of them with
nonzero weight); the negative control that drops the global sign $(-1)^{|\theta|-1}$ disagrees on
$327$ of the $541$.
\end{itemize}
So Theorem~A is the $q$-Murnaghan--Nakayama rule for Iwahori--Hecke algebra characters,
transported by $\omega$. The honest caveat: \textbf{Ram (1991), Halverson (1995) and
Halverson--Ram (1995) were not opened.} None is on arXiv. The earliest source I can name from a
document I have actually read is Jing--Liu \S4.2 (2023), which proves the rule and points at
three earlier ones. Before the word ``classical'' is used in anything outgoing, at least one of
those three should be read at source.

Two nearby results were checked and are \emph{not} it, which is worth recording because both were
on my candidate list: Okada, arXiv \texttt{1904.03386} Cor.~6.8 (attributed there to Morris, JLMS
\textbf{39} (1964) 481--488) and Fayers, arXiv \texttt{2003.07713} Prop.~3.6 both live in the
shifted / Schur $P$--$Q$ world and carry the weight $2^{\#\text{components}}$ with no $t$; they
are the $t=-1$ shadow of the rule, in a different arena. Fayers' remark on MathOverflow~337891
that he holds an unpublished Hall--Littlewood version by duality concerns that shifted rule, not
this one.

\subsection*{Theorem~\ref{thm:main} of this paper}

Not searched. The brief capped the prior-art phase at Theorem~A, and I spent the cap there. What
the browse log of 2026-09-04 records is a \emph{mechanism} for why the commutator across two
different ribbon sizes is unlikely to have been asked: in the Fock-space literature the ribbon
size $n$ is fixed by the choice of $U_q(\widehat{\mathfrak{sl}}_n)$, and the Heisenberg index
counts how many ribbons rather than what size, so $R_e$ and $R_{e'}$ for $e\ne e'$ act on
different Fock spaces there. That is an argument, not a search. The status of
Theorem~\ref{thm:main} is \textbf{unsearched}, and it should not be described as new.

\section{Computational verification}\label{sec:comp}

All code is in \texttt{probes/2026-09-04-Q76/}. The engine \texttt{abacus.py} implements
Corollary~\ref{cor:Rmaya} on beta-numbers; it is checked against \texttt{symfunc.py} from
\texttt{probes/2026-09-03-Q75/}, which builds $R_e$ by direct enumeration of skew shapes on Young
diagrams and shares no code with it. The commutator is then formed by composing operators, with
no reference to Theorem~\ref{thm:main}.

\begin{center}
\begin{tabular}{llr}
\toprule
Claim & Range & Result\\
\midrule
Prop.~\ref{prop:hop}(i)--(ii): beads vs.\ skew shapes, entrywise
  & $|\lambda|\le8$, $e=2..6$ & $335/335$ configs\\
Theorem~\ref{thm:main}, entrywise vs.\ brute force
  & $|\lambda|\le8$, all $2\le e<e'\le7$ & $238594/238594$\\
\quad of which nonzero & same & $11360$\\
\quad two-bead (crossing) entries, shape $\pm t^a(1-t^2)$ & same & $8408$\\
\quad ribbon entries, shape $\pm t^a(1+t)$ & same & $2952$\\
Prop.~\ref{prop:hop}(iv): the cut reading, on Young diagrams
  & $|\lambda|\le6$, $2\le e<e'\le5$ & $2552/2552$\\
Prop.~\ref{prop:hop}(iii): the diagonal interval
  & $|\lambda|\le7$, $e=2..6$ & $1058/1058$\\
Lemma~\ref{lem:wit}, both families & all $2\le e<e'\le8$ & $21/21$, $21/21$\\
\S\ref{sec:prior} dictionary, generating functions & $e\le5$ & exact\\
\S\ref{sec:prior} dictionary, weights & $541$ skew shapes & $541/541$\\
\bottomrule
\end{tabular}
\end{center}

The counts $8408$ and $2952$ agree with the two entry-shape counts computed on 2026-09-04 by an
entirely different route --- direct border-strip enumeration on Young diagrams, over the same
range --- and recorded in the browse log before this session began. That is a cross-check on the
support and the shapes, not on the formula.

\paragraph{Negative controls, each with its variation named before it was run.}
\begin{itemize}
\item \emph{C1, the trivial kernel.} $e=e'$ must give the zero matrix: $120/120$ configurations
($|\lambda|\le6$, $e\le5$). Theorem~\ref{thm:main} does not apply there, and
Lemma~\ref{lem:assign} is exactly where $e\ne e'$ enters.
\item \emph{C2, antisymmetry.} Exchanging $e$ and $e'$ negates every entry: $1208/1208$.
\item \emph{C3, the height convention.} This is the control that tests the part of the object
that is mine. Replacing $h$ by the number of rows (i.e.\ $h+1$) in the bead operator must
\emph{break} Theorem~\ref{thm:main}: it disagrees on $761$ of $6665$ entries. The check is
therefore not blind to the grading. No convention was chosen during this session; the height is
the one fixed by Definition~\ref{def:maya} and the previous paper.
\item \emph{C4, the crossing condition is load-bearing.} A variant prediction that keeps
everything else but drops the crossing hypothesis in case (ii) --- i.e.\ returns
$\pm t^{H-1}(1-t^2)$ for every two-bead pair --- disagrees with the truth on $12088$ entries,
exactly the number of reachable non-crossing pairs. So the vanishing in case (ii) is a real
constraint and not an artefact of an empty support.
\item \emph{C5, the cut condition is load-bearing.} A variant that returns $t^{h-1}(1+t)$ on
every $(e+e')$-ribbon, ignoring $\epsilon_e-\epsilon_{e'}$, disagrees on $7040$ entries.
\item \emph{C6, the prior-art dictionary.} Dropping the global sign $(-1)^{|\theta|-1}$
disagrees on $327/541$ shapes.
\end{itemize}
Every control fires. The one I would treat as decisive is C3: the two shapes $\pm t^a(1+t)$ and
$\pm t^a(1-t^2)$ are statements about the height grading, and a check that survived a perturbed
height would be measuring only the support.

\section{Gaps}

\begin{enumerate}
\item The prior-art status of Theorem~\ref{thm:main} is \textbf{unsearched} (\S\ref{sec:prior}).
\item The identification of the previous paper's Theorem~A as classical rests on Jing--Liu
\texttt{2310.15730} \S4.2, read at source by me, plus a dictionary I verified computationally.
The onward attribution to Ram (1991), Halverson (1995) and Halverson--Ram (1995) is a
citation-chain claim read inside Jing--Liu; none of those three has been opened, and none is on
arXiv.
\item The Hall--Littlewood identification (\texttt{cor:HL}) of the previous paper (the identification $f_e(t)=P_{(e)}(x;-t)$)
still rests on two facts quoted from Macdonald III from memory. The prior-art phase did not reach
them; the partial findings are that Jing--Liu \S5.2 gives the exponential form
$\sum_nq_n(X;t)z^n=\exp\bigl(\sum_n\frac{1-t^n}{n}p_nz^n\bigr)$ and the identification
$q_n=Q_{(n)}$, but not the product form nor $q_r=(1-t)P_{(r)}$; and that their \S4.2 prose
describes $\tilde h_r$ as a one-row Hall--Littlewood function, which corroborates the
identification without supplying the two Macdonald displays.
\item Theorem~\ref{thm:main} says nothing about the $q$-deformed rigid factor $(1+qt)$ of Q59,
nor about the Heisenberg scalar $[e]_{t^2}$ of $[R_e^\perp,R_e]$. Those brackets do not vanish at
$t=-1$; this one does. They are different phenomena and this paper does not connect them.
\item The two-bead entries are described by a condition on beads. Corollary~\ref{cor:diag}
translates it to diagonals, but I have not characterised, intrinsically in terms of the skew shape
$\mu/\lambda$ alone, which shapes of size $e+e'$ arise in case (ii) and with which multiplicities.
For the ribbon case this is done (Corollary~\ref{cor:ribbon}).
\end{enumerate}

\end{document}
